\section{Framed knots}\label{sec:framing}

\subsection{Definition}
% [F1]
\begin{definition}\label{def:framed}
For a partition $\lambda$ the refined framing factor is
\[
\theta_\lambda(q,t)=\bigl(-(q/t)^{1/2}\bigr)^{|\lambda|}\,q^{n(\lambda')}\,t^{-n(\lambda)}
=(-1)^{|\lambda|}\prod_{s\in\lambda}q^{a(s)+\frac12}\,t^{-l(s)-\frac12} .
\]
The refined Chern--Simons invariant at framing $x\in\Z$ is $\brCS^{(x)}_\lambda=\theta_\lambda^x\,\brCS_\lambda$, and likewise $\rCS^{(x)}_\lambda=\theta_\lambda^x\,\rCS_\lambda$.
The quantities obtained by inserting $\brCS^{(x)}_\lambda$ into \eqref{eq:GT1} and \eqref{eq:GT2} are denoted by
$f^{q,(x)}_\mu$, $f^{\bar t,(x)}_\mu$, $\hf^{(x)}_\rho$, $\hbN^{(x)}_{\rho;i,j,k}$ and $\hbN^{(x)}_{\rho,g,\beta,F}$.
The case $x=0$ is that of KN.
\end{definition}

The factor $\theta_\lambda$ is a monomial with coefficient $\pm1$.
Under the Adams operation $\brCS_\lambda\mapsto\brCS_\lambda(a^k,q^k,t^k)$ inside the plethystic logarithm it is replaced by $\theta_\lambda(q^k,t^k)^x$; the sign $(-1)^{|\lambda|}$ is \emph{not} raised to the power $k$.
This is the same treatment as in the example
$f^{(p)}_{\yng(2)}=q^pW_{\yng(2)}-\tfrac12\bigl(W_{\yng(1)}^2+(-1)^pW_{\yng(1)}(q^2,a^2)\bigr)$ of Mari\~no and Vafa \cite{Marino01}
(we write $a$ for their variable $\lambda$).

\subsection{Origin of the framing factor}
% [F2]
\begin{proposition}\label{prop:MV}
At $t=q$ one has $\theta_\lambda^x=(-1)^{|\lambda|x}q^{x\kappa_\lambda/2}$, where $\kappa_\lambda=|\lambda|+\sum_i(\lambda_i^2-2i\lambda_i)=2\sum_{s\in\lambda}c(s)$.
Hence the limit $t=q$ of $\brCS^{(x)}_\lambda$ is the framed invariant $W^{(p)}_R=(-1)^{|R|p}q^{p\kappa_R/2}W_R$ of Mari\~no and Vafa \cite{Marino01} with $p=x$.
\end{proposition}
\begin{proof}
At $t=q$ we have $(q/t)^{1/2}=1$ and $n(\lambda')-n(\lambda)=\sum_s c(s)=\kappa_\lambda/2$.
The direction of the variable ($q$ versus $q^{-1}$) is fixed by comparing the principal specialisation for the unknot with the quantum dimension of Mari\~no and Vafa:
they agree up to a monomial with $q$ identified with $q$, and do not agree with $q\mapsto q^{-1}$ when $\lambda\neq\lambda^T$.
\end{proof}

\begin{remark}[Refined topological vertex and the $T$-matrix]\label{rem:Taki}
The framing factor of the refined topological vertex in the form of Taki \cite{Taki07} is
$f_\mu(t,q)=(-1)^{|\mu|}t^{\|\mu^T\|^2/2}q^{-\|\mu\|^2/2}$.
Using $\|\mu\|^2=2n(\mu')+|\mu|$ one finds $f_\mu(t,q)=\theta_\mu(q,t)^{-1}$:
the factor of Taki is exactly the inverse of that of Definition~\ref{def:framed}, the only difference being the choice of the positive direction of framing.
Taki showed that multiplying the factor $(-1)^{|\nu|}t^{n(\nu)}q^{-n(\nu^T)}$ of Iqbal, Koz\c{c}az and Vafa \cite{Iqbal07} by $(t/q)^{|\nu|/2}$ is necessary for agreement with the K-theoretic Nekrasov partition function.
On the Chern--Simons side, the eigenvalues $T_i=c_T\,q^{(\lambda_{R_i}+\beta\rho)^2/2}$ ($t=q^\beta$) of the refined $T$-matrix of Aganagic and Shakirov satisfy,
up to the $SU(N)$ correction, $T_\lambda/T_\emptyset=(a^{1/2}(q/t)^{1/2})^{|\lambda|}\,q^{n(\lambda')}t^{-n(\lambda)}$.
Thus, apart from a monomial per box, one unit of framing is the action of the refined $T$-matrix.
We use the symbol $\theta_\lambda$ rather than Taki's $f_\mu$ to avoid confusion with $f^q_\mu$, $f^{\bar t}_\mu$ and $\hf_\rho$.
\end{remark}

% [F6]
\begin{proposition}\label{prop:bold}
In the bold variables,
\[
\theta_\lambda=\bt^{\,|\lambda|+2n(\lambda')}\,\bq^{\,2(n(\lambda')-n(\lambda))},
\]
which contains no sign. For example $\theta_{\yng(1)}=\bt$, $\theta_{\yng(2)}=\bq^2\bt^4$, $\theta_{\yng(1,1)}=\bq^{-2}\bt^2$ and $\theta_{\yng(2,1)}=\bt^5$.
\end{proposition}
\begin{proof}
By \eqref{eq:cov}, $-(q/t)^{1/2}=\bt$, $q=\bq^2\bt^2$ and $t=\bq^2$.
\end{proof}
Proposition~\ref{prop:bold} shows that a change of framing shifts the degrees of (coloured) HOMFLY-PT homology without signs.
In the unrefined limit $\bt=-1$ one gets $\theta_\lambda^x\to(-1)^{|\lambda|x}$:
the Mari\~no--Vafa sign is the shadow of the refined degree shift $\bt^{|\lambda|x}$.

\subsection{Structure in the framing}
% [F5]
\begin{lemma}\label{lem:adams}
Let $c$ be a monomial with coefficient $1$ on which the Adams operations act by $c\mapsto c^k$ (for instance $c=a^{1/2}$ or $(q/t)^{1/2}$).
Replacing $\brCS_\lambda$ by $c^{|\lambda|}\brCS_\lambda$ for all $\lambda$ replaces $\hf_\rho$ by $c^{|\rho|}\hf_\rho$.
\end{lemma}
\begin{proof}
The factor $c^{|\lambda|}$ amounts to the substitution $u\mapsto cu$ of the symmetric-function variables, which commutes with the plethystic logarithm because $c$ is compatible with the Adams operations.
\end{proof}
For $c=(q/t)^{1/2}=-\bt$ this shifts the $\bt$-degree of both forms uniformly and changes the overall sign, so it does not affect positivity.

% [F7]
\begin{proposition}\label{prop:expsum}
Let $|\rho|=n$. Then $\bt^{-nx}\hf^{(x)}_\rho$ is an exponential polynomial in $x$, that is, a finite sum $\sum_i c_i\,r_i^{\,x}$ with $c_i$ independent of $x$.
The bases $r_i$ are $\{\bq^2\bt^2,\bq^{-2},1,-1\}$ for $n=2$ and $\{\bq^6\bt^6,\bq^{-6},\bq^2\bt^2,\bq^{-2},\bt^2,1,-1\}$ for $n=3$.
After separating even and odd $x$, the bases are monomials with coefficient $1$.
\end{proposition}
\begin{proof}
$\hf_\rho$ is a linear combination, with fixed rational functions as coefficients, of products of the $\brCS_\lambda$ and their Adams images.
With framing, each term acquires the factor $\prod_i\theta_{\lambda_i}(q^{k_i},t^{k_i})^x$.
For $n=2$ the factors are $\theta_{\yng(2)}=\bq^2\bt^4$, $\theta_{\yng(1,1)}=\bq^{-2}\bt^2$, $\theta_{\yng(1)}^2=\bt^2$ and $\theta_{\yng(1)}(q^2,t^2)=-\bt^2$;
multiplying by $\bt^{-2x}$ gives the stated bases. The case $n=3$ is the same.
\end{proof}
In terms of $v=\bq^2\bt$ the bases $\bq^2\bt^2=v\bt$ and $\bq^{-2}=\bt/v$ form a pair exchanged by $\sigma$.
The dependence on $x$ is of the same type as the dependence of twist knots on the number of twists (\S\ref{sec:two}): framing behaves like one more twist variable.

% [F8, F9]
\begin{proposition}\label{prop:x-indep}
For every $x\in\Z$ the following hold:
{\rm(i)} $\theta_\lambda(1/t,1/q)=\theta_{\lambda^T}(q,t)$, and consequently the transposition symmetry of $\rCS_\lambda$, the equality of the $\hf$ obtained from \eqref{eq:GT1} and from \eqref{eq:GT2},
and the $\sigma$-invariance of $\hf_\rho$ hold for $\hf^{(x)}$ if and only if they hold for $\hf^{(0)}$;
{\rm(ii)} for $|\rho|=2$, $\hf^{(x)}_\rho$ is a Laurent polynomial if and only if $\hf^{(0)}_\rho$ is.
\end{proposition}
\begin{proof}
Statement (i) follows from the definition.
For (ii), note that for $|\rho|=2$, $\hf_\rho$ is a linear combination of the four inputs $\rCS_{\yng(2)}$, $\rCS_{\yng(1,1)}$, $\rCS_{\yng(1)}^2$ and $\rCS_{\yng(1)}(a^2,q^2,t^2)$
divided by the common denominator $2A^2(Q^4-1)(T^4-1)(Q^2T^2-1)$, where $(A,Q,T)=(a,q,t)^{1/2}$.
On the zero locus of each irreducible factor of the denominator exactly two of the inputs survive in the numerator, and their framing factors coincide on that locus (a symbolic computation).
Hence the vanishing of the numerator on the locus does not depend on $x$.
\end{proof}

% [F10]
\begin{proposition}[Necessity of the sign]\label{prop:sign}
Suppose that the sign $(-1)^{|\lambda|x}$ in Definition~\ref{def:framed} is omitted. If $x$ is odd, then on the loci $Q=i$ and $T=i$ (that is, $q=-1$ and $t=-1$)
the framing factors of the two surviving inputs differ, and for generic input $\hf_{\yng(2)}$ and $\hf_{\yng(1,1)}$ are not Laurent polynomials.
For even $x$ the sign has no effect. Omitting the factor $(q/t)^{1/2}$ per box preserves the Laurent property.
\end{proposition}
\begin{proof}
Omitting the sign changes only the factor of the Adams image $\rCS_{\yng(1)}(a^2,q^2,t^2)$, by $(-1)^x$; the coincidence fails on the loci where this term survives.
The last statement is Lemma~\ref{lem:adams}.
\end{proof}
\begin{compresult}
For the trefoil at $x=1$ without the sign, the factor $(1+q)(1+t)$ remains in the denominators of $\hf_{\yng(2)}$ and $\hf_{\yng(1,1)}$.
At $x=2$, and at $x=1$ with the sign, they are Laurent polynomials.
\end{compresult}

Proposition~\ref{prop:sign} is the refined version of the observation of Mari\~no and Vafa \cite{Marino01} that the sign is essential for integrality.
They explained the sign by a change of the flat open-string coordinate, $\hat u\mapsto\hat u+i\pi p$, i.e.\ $V\mapsto(-1)^pV$,
which re-expands the generating function $\sum_RW_R\,\mathrm{Tr}_RV$ and multiplies $W_R$ by the constant sign $(-1)^{|R|p}$.
A constant sign is not raised to the $d$-th power in the Adams image $W_R(q^d,a^d)$, so this is not a trivial replacement of the type of Lemma~\ref{lem:adams}:
on the BPS side the $d$-fold cover terms acquire $(-1)^{d|R|}$. The sign in Definition~\ref{def:framed} is treated in the same way.
The remaining part $(q/t)^{|\lambda|/2}q^{n(\lambda')}t^{-n(\lambda)}$ of $\theta_\lambda$ comes from the eigenvalues of the refined $T$-matrix (Remark~\ref{rem:Taki}), which contain no sign.
Aganagic and Shakirov \cite{Aganagic11} remark that the prefactor $\sqrt{q/t}$ for the trefoil should also be related to a change of flat coordinates.
We do not derive the refined change of flat coordinates from geometry.

\subsection{A criterion for all framings}\label{sec:newton}
Proposition~\ref{prop:expsum} reduces positivity for all $x$, for a fixed knot, to finitely many checks.
Split $x$ by parity and direction ($x\ge x_0$ or $x\le x_0$), write $x=x_0\pm2y$ with $y\ge0$, order the bases of the resulting exponential polynomial in $y$ as $(r_0,r_1,\dots)$ (they are the monomials $r_i^{\pm2}$ of Proposition~\ref{prop:expsum}) and form the Newton (divided-difference) expansion
\[
\bt^{-nx}\hf^{(x)}_\rho=\sum_{k}c_k\,h_{y-k}(r_0,\dots,r_k),
\]
where $h_j$ is the complete homogeneous symmetric polynomial ($h_j=0$ for $j<0$).
If all $c_k$ are Laurent polynomials whose coefficients have one common sign, then, since the $h_j$ of monomials with coefficient $1$ have non-negative coefficients,
the spin form is non-negative for all $y\ge0$.
$\sigma$-invariance for all $x$ follows from Proposition~\ref{prop:x-indep}, and Lemma~\ref{lem:G} then gives non-negativity of the genus form.
The finitely many $x$ between the starting point $x_0$ and $0$ are checked directly.
We refer to this as the \emph{Newton criterion}; it is used in \S\ref{sec:two} both for the framing and for the number of twists.

\subsection{Physical picture}\label{sec:physics}
The refined large $N$ duality for framed knots can be assembled from two ingredients.
On the Chern--Simons side, $\theta_\lambda$ without its sign is the action of the $T$-matrix of refined Chern--Simons theory (Remark~\ref{rem:Taki}).
In the construction of Aganagic and Shakirov the $U(1)_R$ symmetry exists because the three-manifold admits a semi-free $U(1)$ action (it is a Seifert manifold),
and a change of framing of a torus knot stays within the same formalism of $S$ and $T$ matrices.
On the string side, in the argument of Mari\~no and Vafa the framing integer corresponds to the choice of the flat coordinate of $L_K$ (the boundary condition at infinity), and the sign $(-1)^{|\lambda|x}$ comes from this change of coordinate.

Two points are not derived in this paper:
that the Schwinger computation of KN and the extraction formulae \eqref{eq:GT1}, \eqref{eq:GT2} apply unchanged for all $x$,
and that the action of $U(1)_R$ on the string side remains trivial.
Under these assumptions positivity is expected for all framings. The precise conjectures are stated in \S\ref{sec:conj} after the evidence has been presented.
